\chapter{ماذا تفعل الآلة في الليل}
% full version: chapters/20_sleep.tex

النوم ليس إطفاء للآلة، بل تغيير للنمط. ويُرى هذا في المحطات الإذاعية. في النوم العميق تُخفض كلها تقريبًا: الأسيتيل كولين والنورأدرينالين والسيروتونين. وفي نوم الأحلام يعود الأسيتيل كولين إلى مستواه النهاري، ويكاد النورأدرينالين والسيروتونين يصمتان، ويُعطَّل المخرج إلى العضلات بالمكابح، ولهذا لا نركض في النوم خلف ما نهرب منه في الحلم. كل هذا مقيس.

\section*{متى يحين وقت النوم}

تخيّل مكثفًا يُشحن طوال النهار ما دمت مستيقظًا، ويُفرغ في الليل. الشحنة هي ”ضغط النوم“. تنام حين تبلغ الشحنة العتبة العليا، وتستيقظ حين تهبط إلى الدنيا. أما العتبتان نفساهما فتصعدان وتهبطان بسلاسة مع الإيقاع اليومي من الفصل 16. ومخطط بسيط كهذا يصل وحده إلى ست عشرة ساعة من اليقظة وثماني ساعات من النوم. ويفسر لماذا ننام بعد ليلة بلا نوم لا ضعف المدة بل أعمق: المكثف الممتلئ يُفرغ أسرع. والمرشح لدور ”الشحنة“ مادة تتراكم في الدماغ في أثناء العمل؛ أما أنها هي بالذات فلا يزال فرضية.

\pencilfig{120}{%
% figures/sleep_{S,H,L}.dat from models/sleep_models.py, t in hours
\begin{scope}[x=0.1cm, y=2.6cm]
\foreach \a/\b in {0/4.3,21.2/29.3,45.4/53.4,69.5/72}{\fill[pcBlue, opacity=0.1] (\a,0) rectangle (\b,1);}
\draw[pen] (0,0) -- (72,0);
\draw[penlite=pcRed, densely dashed] plot file {figures/sleep_H.dat};
\draw[penlite=pcGreen, densely dashed] plot file {figures/sleep_L.dat};
\draw[pcBlue, line width=1.0pt] plot file {figures/sleep_S.dat};
\foreach \t/\l in {0/0,24/يوم,48/يومان,72/3 أيام}{\draw[pcGray, line width=0.5pt] (\t,0) -- (\t,-0.04); \node[lbl, anchor=base] at (\t,-0.16) {\l};}
\node[lbl, text=pcBlue] at (25.25,0.93) {نوم}; \node[lbl, text=pcBlue] at (49.4,0.93) {نوم};
\end{scope}
\node[lbl, text=pcRed, anchor=west] at (7.3,2.0) {النعاس};
\node[lbl, text=pcGreen!70!black, anchor=west] at (7.3,0.62) {الاستيقاظ};
\node[lbl, text=pcBlue, anchor=west] at (7.3,1.35) {ضغط النوم};
}{ضغط النوم يتراكم في النهار ويذوب في الليل؛ ننام عند العتبة العليا، ونستيقظ عند الدنيا، والعتبتان تتأرجحان مع اليوم.}

\section*{الأرجوحة البطيئة}

في النوم العميق تشتغل مناطق كاملة من الدماغ وتنطفئ بتزامن، أقل من مرة في الثانية. وهذا مخطط مألوف: مجموعة تثير نفسها وتتعب هي مذبذب. في النهار يكبت الأسيتيل كولين تعب العصبونات، فتعمل المجموعة نفسها بانتظام؛ وفي الليل يقل، فتبدأ الشبكة بالتأرجح.

\pencilfig{20}{\pencilart{20}}{في الليل تتأرجح الشبكة ببطء: تارة تعمل كل العصبونات تقريبًا، وتارة تصمت كلها تقريبًا.}

\section*{التسريع إلى الأمام}

في هذه الأرجوحة يجري أمتع ما في الأمر. سلاسل تنشيط العصبونات التي وقعت في النهار ”تُعاد“ من جديد، وأسرع بأضعاف: في منطقة الذاكرة نحو عشرين مرة، وفي القشرة بضع مرات. وإذا عُزز اصطناعيًا توافق هذه الإعادات مع الأرجوحة البطيئة، تحسنت الذاكرة: وهذه تجربة سببية. لماذا يلزم هذا؟ فرضية معقولة: الذاكرة الطويلة الأمد البطيئة تتعلم من الإعادات ممزوجة بالقديم، كي لا يمحو الجديدُ السابق. ويعرف المهندسون هذه الآفة: الشبكة العصبية المدرَّبة على مهمة جديدة تنسى القديمة إذا لم تُكرَّر الأمثلة القديمة.

\section*{تقليم الوصلات}

وأمر آخر مقيس: بعد النوم تصغر نقاط التماس بين العصبونات في المتوسط، بنسبة من بضعة في المئة إلى عشرات في المئة، بالتناسب مع حجمها. كأن كل مستويات الصوت خُفضت قليلًا بحركة واحدة: بقيت النسب، وبقي ما تُعلّم، وتحرر مكان لتعلم الغد. أما أن هذه هي المهمة الرئيسية للنوم ففرضية.

\section*{الأحلام والتنظيف}

نوم الأحلام يشبه اختبارًا على منصة التجربة: الآلة تعمل بكامل طاقتها، لكن المداخل مستبدلة بإشارات داخلية، والمخرج معطَّل. لماذا يلزم هذا؟ لا نعرف. وفكرة أن الدماغ ”يُغسل“ من الفضلات في النوم لا تزال موضع جدل: هناك تجارب معها وأخرى ضدها.

وأعجب ما في الأمر: النوم قد يكون موضعيًا. بعد يقظة طويلة، ”تغفو“ مناطق منفردة من قشرة الإنسان المستيقظ لحظة، فتصمت في منتصف العمل.

\fullversion{20}
